\subsection*{A -- Premier placement.}
	\begin{enumerate}
		\item $a_{1} = a_0+200 = \np{20200}$ \quad et \quad $a_{2} = a_1 +200 = \np{20400}$.

		\item \begin{enumerate}
			\item Chaque année, le capital augmente de 200\euro{}, on a donc pour tout entier naturel $n$ :

			$a_{n+1}=a_n+200$.

			\item On reconnaît la relation de récurrence d'une suite arithmétique. La suite $(a_n)$ est arithmétique de raison $r=200$.
		\end{enumerate}

		\item La suite $(a_n)$ est arithmétique donc on a pour tout $n$ entier naturel :\quad $a_n=a_0+200n$.

		Soit ici,\quad $a_n=\np{20000} +200n$.

		\item On cherche $n$ tel que $a_n \geqslant  \np{22000}$.\hspace*{2cm}
		$\begin{aligned}[t] a_n \geqslant  \np{22000} & \iff \np{20000} +200n \geqslant  \np{22000}\\
		 & \iff 200n \geqslant  \np{2000}\\
		 & \iff n \geqslant \dfrac{\np{2000}}{200} \vspace*{2mm}\\
		 & \iff n \geqslant  10
		\end{aligned}$

		Ils devront attendre 10 ans soit à partir de 2035.
	\end{enumerate}

\subsection*{B -- Second placement.}
	\begin{enumerate}
		\item Le capital augmente de $2\%$, il est donc multiplié par 1,02.

			$b_1=b_0 \times 1,02 = \np{20000} \times 1,02=\np{20400}$.

		\item \begin{enumerate}
			\item En généralisant l'opération, chaque année, le capital est obtenu en multipliant le capital de l'année précédente par 1,02 :

			Pour tout $n$ entier naturel :\quad $b_{n+1}=1,02 b_n$.

			\item On reconnaît la relation de récurrence d'une suite géométrique.

			La suite $(b_n)$ est géométrique de raison $q=1,02$.
		\end{enumerate}

		\item La suite $(b_n)$ est géométrique donc on a, pour tout entier naturel $n$: \quad
		$b_n=b_0 \times q^n$

		soit $b_n= \np{20000} \times 1,02^n$

		\item Le tableau permet de dire que dès la 5\up{ème} année, ils auront le capital nécessaire à leur projet, soit en 2030.
	\end{enumerate}

\subsection*{C -- Bilan}
	Le deuxième placement est plus intéressant pour eux.
